\subsection{ソフトウェア工学方法}
ソフトウェア工学方法は、対象コンピュータ向けのソフトウェアを開発するための組織化された体系的な進め方である。方法の選択は、プロジェクトの成功に大きく影響する。

SWEBOK 第11章では、代表的な方法を、経験則的方法、形式手法、プロトタイピング方法、アジャイル方法に分けて説明している。

\subsubsection{経験則的方法}
経験則的方法とは、実務経験に基づいて広く使われている方法である。ここでいう経験則は、単なる勘ではない。過去の開発で有効だった表現、分解、設計、確認のやり方を体系化したものである。

\paragraph{構造化分析・設計方法}
構造化分析・設計方法は、主に機能や振る舞いの視点からモデルを作る。高水準のソフトウェアモデルから始め、データ要素や制御要素を含む構成要素を段階的に分解していく。詳細化が進むと、手作業または自動生成でコード化し、ビルドし、テストし、検証できる仕様に近づく。

\paragraph{データモデリング方法}
データモデリング方法は、使われるデータや情報を中心にモデルを作る。データ表、実体、属性、関係を定義し、データ要求やデータベース設計を支援する。業務ソフトウェアでは、データが事業上の資源や資産として管理されるため、データモデリングは特に重要である。

\paragraph{オブジェクト指向分析・設計方法}
オブジェクト指向分析・設計方法は、データと振る舞いをまとめたオブジェクトの集合としてソフトウェアを考える。オブジェクトは、現実世界の対象でも、仮想的な対象でもよい。オブジェクトはメソッドを通じて相互作用する。

モデルは複数の図で構成され、段階的に詳細設計へ洗練される。詳細設計は、反復を通じて発展する場合もあれば、実装ビューへ変換される場合もある。代表的な方法には統一プロセス（Unified Process）や、その具体化であるRational Unified Processがある。

\paragraph{アスペクト指向開発方法}
アスペクト指向開発は、横断的関心事を分離して扱う方法である。横断的関心事とは、ログ記録、認証、トランザクション、例外処理、監視など、複数の部品にまたがって現れる関心事である。

横断的関心事を通常の部品の中に散らばらせると、関心事の散乱と絡まりが起こる。アスペクトは、こうした横断的関心事をカプセル化する単位である。ソフトウェアレベルでは、ウィーバが、指定された結合点に助言を織り込む。

\paragraph{モデル駆動開発・モデルベース開発方法}
モデル駆動開発（MDD）は、モデルを開発プロセスの主要成果物として使う方法である。実装や他のモデルは、モデルから半自動または自動で変換されることがある。

モデルベース開発（MBD）は、モデルを用いてシステムを分析する方法である。モデルが必ずしも主要成果物とは限らない。制御、信号処理、通信などの動的システムでは、実行可能仕様やシミュレーションに焦点を当てることがある。モデルからテストケースを生成する場合もある。

\par\smallskip
\begin{center}
\centering
\IfFileExists{assets/generated_figures/ch11/fig-ch11-07.png}{\includegraphics[width=0.96\linewidth]{assets/generated_figures/ch11/fig-ch11-07.png}}{\fbox{\parbox[c][48mm][c]{0.9\linewidth}{\centering 画像生成待ち\\経験則的方法の分類図}}}
\par\smallskip
\refstepcounter{figure}\label{fig:ch11-07}図 \thefigure: 経験則的方法の分類図
\end{center}
図 \thefigure は、経験則的方法の分類図を視覚的に整理したものである。
\par\smallskip

\subsubsection{形式手法}
\term{形式手法}{Formal Method}とは、厳密な数学的記法と言語を使って、ソフトウェアを仕様化し、開発し、検証する方法である。仕様言語を使うことで、モデルの整合性、完全性、正しさを自動または半自動で確認できる場合がある。

形式手法は、曖昧さを減らしたい場面、安全性やセキュリティが重要な場面、通常のテストだけでは十分な確信を得にくい場面で価値が高い。一方で、数学的な表現を理解する負担があるため、すべてのプロジェクトに同じ程度で適用すればよいわけではない。

\paragraph{仕様言語}
仕様言語は、形式手法の数学的基盤を提供する。仕様言語は、要求分析や設計の段階で、入出力の振る舞いを厳密に表すために使われる。通常の第3世代プログラミング言語とは異なり、直接実行するための言語ではないことが多い。記法、構文、意味論、許される関係の集合を含む。

\paragraph{プログラム洗練と導出}
プログラム洗練とは、高水準の仕様を、変換の列によって低水準で詳細な仕様へ近づけることである。変換を重ねることで、最終的に実行可能な表現へ導く。これが可能になるのは、仕様の意味が精密に定義されているからである。

\paragraph{形式検証}
形式検証は、モデルが特定の性質を持つことを数学的に確かめる方法である。モデル検査は形式検証の一例であり、状態空間探索や到達可能性分析を用いる。たとえば、イベントやメッセージの到着順がどのように入れ替わっても、正しい振る舞いを保つかを確認する。

\paragraph{論理推論}
論理推論は、重要な設計ブロックの周りに事前条件と事後条件を定め、すべての入力に対してそれらが成り立つことを数学的論理で示す方法である。これにより、ソフトウェアを実行しなくても振る舞いを予測できる。

\paragraph{軽量形式手法}
軽量形式手法は、実用性と厳密な検証のバランスを取る方法である。たとえば Alloy は、精密で表現力のある記法を用いつつ、定理証明ではなく自動分析によってすばやいフィードバックを与える。ただし、有限の範囲だけを調べるため、完全な証明ではない。

\par\smallskip
\begin{center}
\centering
\IfFileExists{assets/generated_figures/ch11/fig-ch11-08.png}{\includegraphics[width=0.96\linewidth]{assets/generated_figures/ch11/fig-ch11-08.png}}{\fbox{\parbox[c][48mm][c]{0.9\linewidth}{\centering 画像生成待ち\\形式手法の流れを示す図}}}
\par\smallskip
\refstepcounter{figure}\label{fig:ch11-08}図 \thefigure: 形式手法の流れを示す図
\end{center}
図 \thefigure は、形式手法の流れを示す図を視覚的に整理したものである。
\par\smallskip

\subsubsection{プロトタイピング方法}
プロトタイピングとは、不完全または最小限の機能を持つソフトウェアの試作品を作り、理解や確認に使う活動である。目的は、新機能の試行、要求や利用者インタフェースへのフィードバック獲得、設計や実装選択肢の探索などである。

プロトタイピングの特徴は、最もよく理解されていない部分から先に扱う点である。通常の開発方法では、よく理解された部分から始めることが多い。プロトタイプは、広範な再作業やリファクタリングなしに最終製品になるとは限らない。

\par\smallskip\noindent\refstepcounter{table}\label{tab:ch11-07}表 \thetable: プロトタイピング方法の整理\par\smallskip
表 \thetable は、プロトタイピング方法の整理を比較・確認しやすい形で整理したものである。\par\smallskip
\begin{longtable}{>{\raggedright\arraybackslash}p{0.24\textwidth}>{\raggedright\arraybackslash}p{0.70\textwidth}}
\toprule
観点 & 説明 \\
\midrule
プロトタイピング様式 & 使い捨てコード、紙の試作品、動く設計の発展形、実行可能仕様などがある。必要な結果、品質、緊急度に応じて選ぶ。\\
プロトタイピング対象 & 要求仕様、アーキテクチャ要素、アルゴリズム、人間機械インタフェースなど、確認したい具体的対象である。\\
評価技法 & 要求に照らして試す、実装と比較する、将来の設計モデルとして使うなど、作成目的に応じて評価する。\\
\bottomrule
\end{longtable}

\begin{pointbox}{プロトタイプ利用時の注意}
プロトタイプは理解を早めるが、利用者が「もう完成している」と誤解することがある。試作品の目的、捨てる範囲、最終製品との違いを明確にしておく必要がある。
\end{pointbox}

\par\smallskip
\begin{center}
\centering
\IfFileExists{assets/generated_figures/ch11/fig-ch11-09.png}{\includegraphics[width=0.96\linewidth]{assets/generated_figures/ch11/fig-ch11-09.png}}{\fbox{\parbox[c][48mm][c]{0.9\linewidth}{\centering 画像生成待ち\\プロトタイピングの三観点を示す図}}}
\par\smallskip
\refstepcounter{figure}\label{fig:ch11-09}図 \thefigure: プロトタイピングの三観点を示す図
\end{center}
図 \thefigure は、プロトタイピングの三観点を示す図を視覚的に整理したものである。
\par\smallskip

\subsubsection{アジャイル方法}
\term{アジャイル方法}{Agile Method}は、重厚で計画中心の方法に伴う大きな負担を減らすため、1990年代に発展した。短い反復サイクル、自己組織化チーム、単純な設計、コードのリファクタリング、テスト駆動開発、頻繁な顧客関与、各サイクルで動く製品を示すことを重視する。

アジャイル方法は、Plan-Do-Check-Act（PDCA）の改善サイクルをソフトウェア工学へ適用したものとして見ることもできる。EVOは、PDCAを増分的に実装する早期のアジャイル方法の一つとして説明できる。

\paragraph{代表的なアジャイル方法}
\par\smallskip\noindent\refstepcounter{table}\label{tab:ch11-08}表 \thetable: 代表的なアジャイル方法の整理\par\smallskip
表 \thetable は、代表的なアジャイル方法の整理を比較・確認しやすい形で整理したものである。\par\smallskip
\begin{longtable}{>{\raggedright\arraybackslash}p{0.20\textwidth}>{\raggedright\arraybackslash}p{0.74\textwidth}}
\toprule
方法 & 要点 \\
\midrule
RAD & 主にデータ集約型の業務システムで使われる。専用のデータベース開発ツールによって、業務アプリケーションを素早く開発、テスト、展開する。\\
XP & ストーリーやシナリオで要求を扱い、先にテストを作り、顧客が直接関与し、ペアプログラミング、継続的リファクタリング、継続的統合を行う。\\
Scrum & プロジェクト管理に向いたアジャイル方法である。スプリントは最大30日程度で、プロダクトオーナー、スクラムマスター、プロダクトバックログ、日次スクラムなどを用いる。\\
FDD & モデル駆動で短い反復を行う方法である。全体モデル、機能一覧、機能開発計画、機能ごとの設計、コード化・テスト・統合の五段階を重視する。\\
Lean & トヨタ生産方式に由来するリーン製造の考え方をソフトウェア開発へ適用する。最小実用製品を出し、利用者から学び、全体の価値の流れを最適化する。Kanbanはワークフロー管理と可視化を重視する軽量な方法である。\\
\bottomrule
\end{longtable}

\paragraph{アジャイル方法と計画駆動方法の使い分け}
アジャイル方法が有効な場面もあれば、重厚で計画駆動の方法が必要な場面もある。方法選択は、組織の事業上の必要、規制、安全性、契約、チーム規模、要求の安定性、変更の頻度によって決まる。

近年は、アジャイルと計画駆動方法を組み合わせる方法も増えている。大規模アジャイルや企業全体のアジャイルは、多数のアジャイルチームを調整しつつ、アジャイルの利点を保つことを目指す。

\paragraph{リリースエンジニアリングとDevOps}
アジャイル方法は短いリリースサイクルにつながる。リリースエンジニアリングは、ソースコードをコンパイルし、組み立て、完成した製品や部品として届けることに関わる下位分野である。統合、ビルド、テストが中心になる。

DevOpsは、アジャイルや継続的デプロイと混同されることがある。区別して考えるためには、リリース管理を、開発と運用の協働の隙間を埋める方法として見るとよい。継続的統合システムのような自動化管理ツールは、展開スケジュールに合わせてリリースを管理するうえで重要である。

\par\smallskip
\begin{center}
\centering
\IfFileExists{assets/generated_figures/ch11/fig-ch11-10.png}{\includegraphics[width=0.96\linewidth]{assets/generated_figures/ch11/fig-ch11-10.png}}{\fbox{\parbox[c][48mm][c]{0.9\linewidth}{\centering 画像生成待ち\\アジャイル方法の比較図}}}
\par\smallskip
\refstepcounter{figure}\label{fig:ch11-10}図 \thefigure: アジャイル方法の比較図
\end{center}
図 \thefigure は、アジャイル方法の比較図を視覚的に整理したものである。
\par\smallskip

\par\smallskip
\begin{center}
\centering
\IfFileExists{assets/generated_figures/ch11/fig-ch11-11.png}{\includegraphics[width=0.96\linewidth]{assets/generated_figures/ch11/fig-ch11-11.png}}{\fbox{\parbox[c][48mm][c]{0.9\linewidth}{\centering 画像生成待ち\\アジャイル、リリースエンジニアリング、DevOpsの関係図}}}
\par\smallskip
\refstepcounter{figure}\label{fig:ch11-11}図 \thefigure: アジャイル、リリースエンジニアリング、DevOpsの関係図
\end{center}
図 \thefigure は、アジャイル、リリースエンジニアリング、DevOpsの関係図を視覚的に整理したものである。
\par\smallskip
